% compendium/execution-market-making/order-flow-statistics.tex
% Entries migrated on 2026-10-16 from the legacy stylised-facts, trade-duration-and-event-time and
% statistical-models-for-order-flow files (slice 33). Owners: market-simulation (stylised facts, volatility estimation,
% deseasonalisation, clocks, point processes) and market-making-lob (Hawkes processes for order flow).
% Dispositions in notes/migration-slices.md; the legacy files verbatim in notes/artefacts/market-simulation/ and
% notes/artefacts/market-making-lob/; numerical checks in notes/artefacts/market-simulation/order-flow-checks.py.

\section{Order flow in time: stylised facts, clocks and point processes}

\begin{qfevidence}[Stylised facts of high-frequency prices and order flow: heavy tails, volatility clustering, long memory of order signs, activity clustering and intraday patterns, with their samples]\label{ev:high-frequency-stylised-facts}
\qffield{Statement} Bouchaud et al.\ report, each with its own sample:
\begin{itemize}
\item \emph{heavy tails}: unconditional return distributions decay as a power law; between about a minute and a few hours a Student-$t$ with tail parameter $\mu$ near 3 fits ``reasonably well'', and their figure of normalised daily returns of U.S.\ stocks, CDS spreads and implied volatilities for 2011--2013 shows tails with exponent about 3.5 \citep[chap.~2, sec.~2.2.1, pp.~28--30]{catalogue-0528};
\item \emph{volatility clustering}: returns are nearly uncorrelated but far from IID; writing $r_t=\sigma_t\xi_t$, log-volatility has a variogram growing like $\ln(1+\tau)$ up to a cut-off of years, so bursts occur on all time scales (sec.~2.2.2, pp.~30--31);
\item \emph{long memory of order signs}: the autocorrelation of market-order signs in event time decays roughly as $\ell^{-\gamma}$ with $\gamma<1$, close to $\ell^{-0.5}$ for four NASDAQ stocks in 2015; with $C(\ell)\approx0.5\,\ell^{-1/2}$, $C(10{,}000)\approx0.005$ (chap.~10, sec.~10.1, pp.~187--189);
\item \emph{activity clustering}: event arrivals cluster beyond what intraday patterns explain; counts in windows have variance well above their mean (chap.~9, secs.~9.1.2 and~9.3.3, pp.~165 and~174--175);
\item \emph{intraday patterns}: market-order and limit-order volume follow a J-shaped profile over the day, the busiest periods about four times the quietest, for PCLN, TSLA, CSCO and INTC on NASDAQ in 2015 (chap.~4, sec.~4.2.1, p.~61); volume at the best quotes behaves differently for large- and small-tick stocks (same page).
\end{itemize}

\qffield{Conditions} Long memory of signs and clustering of activity are different facts: Poisson arrivals can carry persistent signs, and uncorrelated signs can arrive in clusters (chap.~10, pp.~187--188). Each fact is reported for a sample and a method; ``universality'' of the tail exponent is Bouchaud et al.'s reading of many studies, whose estimates vary (p.~29). Hasbrouck notes the same intraday seasonality, volume and volatility elevated at the start and end of sessions \citep[note~1 to chap.~4, p.~180]{s6-hasbrouck}.

\qffield{Practice} For a simulator, these are the targets the pathway's outcome names: a reactive simulation is checked against them, not tuned until it matches one of them by construction. They describe data; none is by itself a trading signal. \emph{What simulators are scored on} (read 2026-10-16). Vyetrenko et al.\ catalogue realism checks in five groups: return distributions (absence of linear autocorrelation beyond about 20 minutes, heavy tails with aggregational normality, volatility clustering, absolute-return autocorrelation decaying as $\tau^{-\beta}$ with $\beta\in[0.2,0.4]$, gain/loss asymmetry, volume--volatility and return--volatility correlations); volumes and order flow (best-queue volumes, order sizes, counts per window, inter-arrival times, placement distance and order lifetimes); non-stationary patterns (intraday and seasonal volume, event days); impact (average mid-price move against binned participation $P_\tau=|V_{\text{buy}}-V_{\text{sell}}|/(V_{\text{buy}}+V_{\text{sell}})$, fitted as $M_i\sim\alpha P_i^{\beta}$); and cross-asset correlation \citep[arXiv v1, sec.~2, pp.~4--6]{w2-vyetrenko-realism}. Their comparisons are visual and distributional, without a single score: two ABIDES configurations against minute returns of 50 U.S.\ stocks resampled daily over 2011 and JPM's NASDAQ book in June 2019, where neither configuration reproduced the spike of near-zero return autocorrelations, the volume--volatility correlation or the U-shaped intraday volume (the richer one reversed it) (secs.~3--4, pp.~6--8). Noble, Rosenbaum and Souilmi validate a queue-reactive simulator on the event-type mix, the imbalance before trades, hourly volume, five-minute realised volatility and the five-minute return distribution, calibrated and compared on the same Databento MBP-10 data for PFE (also INTC, VZ and T), December 2023 to December 2025, 10:00--15:30, with parameter stability checked on disjoint six-month windows; the baseline's return tails are too light, since an ergodic queue-reactive mid-price is asymptotically Brownian \citep[arXiv v1 of 25 March 2026, secs.~2.7--2.8, pp.~8--9]{exp-noble-rosenbaum-souilmi-reality-gap}. They then fit inter-event times to the data, finding a mode at about 29~$\mu$s ($\log_{10}\Delta t\approx4.47$ in nanoseconds) with roughly 25 percent of events faster than that, read as latency races, and add an impact state $\phi_t=\sum_{t_k<t}G(t-t_k)\varepsilon_kV_k$ with a power-law kernel that tilts trade probabilities toward reversion (secs.~3.1 and~4.1, pp.~11--13). Their validation is in-sample except for that stability check, and the paper is an unrefereed preprint.

\qffield{Pitfalls} The legacy taxonomy listed ``long memory'' of ``signs, durations and volatility'' together; the passages read support long memory of order signs and of volatility and clustering of activity, each with a different measure. Durations' own persistence was not checked here.

\qffield{See also} \hyperref[meth:realised-variance-signature-plot]{realised variance and the signature plot}; \hyperref[meth:intraday-deseasonalisation]{intraday deseasonalisation}; \hyperref[mod:hawkes-order-flow]{Hawkes processes for order flow}; \hyperref[mod:propagator]{propagator model} (sign persistence and impact); \hyperref[meth:activity-sampled-bars]{activity-sampled bars} (heavy tails under different clocks); \hyperref[def:intraday-overnight-returns]{intraday and overnight returns}.

\qffield{Further reading} Vyetrenko et al.'s realism metrics and Noble, Rosenbaum and Souilmi's reality-gap tests, readings of \texttt{market-simulation}, were read in their arXiv versions (Practice); the ACM proceedings version of the former was not compared. The legacy file is kept verbatim in \nolinkurl{notes/artefacts/market-simulation/}.

\qffield{Index} stylised facts; heavy tails; fat tails; tail exponent; Student-t; volatility clustering; long memory; order-sign autocorrelation; activity clustering; clustering ratio; intraday seasonality; J-shaped profile; U-shaped profile

\qffield{Status} Verification: partly checked 2026-10-16 (against the text layer of a local copy of \texttt{catalogue-0528}, Cambridge 2018, pp.~28--31, 58--61, 165, 174--175 and 187--189, and \texttt{s6-hasbrouck}, 2007 Oxford printing, p.~180; source follow-up of 2026-10-16: Vyetrenko et al., arXiv v1, pp.~4--8, and Noble, Rosenbaum and Souilmi, arXiv v1, pp.~2 and 8--13). Stability: draft. Currency: samples of 2011--2015, 2019 and 2023--2025.
\end{qfevidence}

\begin{qfmethod}[Realised variance and covariance at high frequency: the signature plot, microstructure noise and the Epps effect]\label{meth:realised-variance-signature-plot}
\qffield{Statement} For a log price $X=\log S$ observed over $[0,T]$ with step $\Delta$ and returns $r_\Delta(j)=X_{j\Delta}-X_{(j-1)\Delta}$, Lehalle and Laruelle state that ``in the It\^o semi-martingale framework'' the realised variance $\hat V_R(\Delta)=\sum_{j=1}^{T/\Delta}r_\Delta(j)^2$ converges in probability to the integrated variance $\int_0^T\sigma_s^2\,ds$ as $\Delta\to0$, and the realised covariance $\hat C_R(\Delta)=\sum_jr^1_\Delta(j)r^2_\Delta(j)$ to $\int_0^T\rho_s\sigma^1_s\sigma^2_s\,ds$ \citep[app.~A, secs.~A.13.1--A.13.2, pp.~316--318]{alpha-laruelle}. Interpretation, standard and checked below: those limits hold for continuous It\^o processes; for a semimartingale with jumps $\hat V_R$ converges to the quadratic variation $\int_0^T\sigma_s^2\,ds+\sum_{s\le T}(\Delta X_s)^2$ and $\hat C_R$ to $\int_0^T\rho_s\sigma^1_s\sigma^2_s\,ds+\sum_{s\le T}\Delta X^1_s\Delta X^2_s$, the last term from co-jumps. If observed prices are $X=M+\varepsilon$ with $M$ a semimartingale and $\varepsilon$ noise, $\hat V_R$ adds the noise terms $\sum\eta_\Delta(j)^2$ and a cross term, and on Rio Tinto data (1 January to 19 June 2012) the estimator diverges as $\Delta\to0$ (pp.~316--318). The covariance estimator needs synchronous data, which quotes and trades are not; the \emph{Epps effect} is Epps's (1979) finding that ``correlations among price changes \dots\ are found to decrease with the length of the interval for which the price changes are measured'', shown for Total and France Telecom on Euronext Paris, 23--27 July 2012 (pp.~318--319, Figure~A.26).

\qffield{Notation} Bouchaud et al.'s \emph{volatility signature plot} is $\sigma^2(\tau)=V(\tau)/\tau$ against the lag $\tau$, with $V(\tau)$ the price variogram: flat for uncorrelated returns, rising with trends and falling with mean reversion or with mean-reverting high-frequency noise \citep[chap.~2, secs.~2.1.2--2.1.3, pp.~25--27]{catalogue-0528}. Lehalle and Laruelle plot $\hat V_R(\Delta)$ against $\Delta$; the two plots carry the same information in different units.

\qffield{Conditions} \emph{Jumps.} Checked (\nolinkurl{order-flow-checks.py}): a unit-jump Poisson path with five jumps on $[0,1]$ has zero diffusion variance, yet its realised variance is 9 on a 10-step grid (two jumps in one interval) and 5, the jump count, from $10^4$ steps on; three common jumps of $+1$ and $-0.5$ give a realised covariance of $-1.5$. \emph{Noise.} The finite-grid expectation below is exact for a driftless Brownian $M$ with constant volatility and observation noise i.i.d., centred and independent of $M$; with drift it adds the squared drift increments, with autocorrelated noise (as in Bouchaud et al.'s mean-reverting noise) the noise term is no longer $2n\sigma_\eta^2$, and for a general semimartingale only the limits above are claimed. Explanations offered for the Epps effect are lead--lag between related assets, asynchronous transactions, and tick and other microstructure effects (Lehalle and Laruelle, p.~318). Checked (\nolinkurl{order-flow-checks.py}): in the simulated model, $X=M+\varepsilon$ with $M$ a driftless Brownian motion of volatility $\sigma$ and $\varepsilon$ i.i.d.\ centred noise of standard deviation $\sigma_\eta$ independent of $M$, $E[\hat V_R]=\sigma^2T+2n\sigma_\eta^2$ exactly for $n=T/\Delta$ returns (each noise difference has variance $2\sigma_\eta^2$ and the cross terms have mean zero), and simulated values track it from 13 to 23{,}400 returns a day; two Brownian motions with correlation 0.6, observed at independent random times and sampled at the previous tick, show a mean correlation over 40 paths of 0.12 at a step of 5 time units against 0.63 at 600, from asynchrony alone.

\qffield{Practice} Bouchaud et al.'s E-mini S\&P~500 signature plot is almost flat from one second to about 20 trading days, decreasing by about 20 percent, which indicates weak mean reversion (sec.~2.1.4, pp.~27--28): for that contract, noise inflates short-horizon volatility only modestly. Whether the plot rises at fine steps depends on the price used (trades bounce between bid and ask; mid-quotes much less) and on the instrument. Lehalle and Laruelle list range-based, multi-scale and bid--ask-model estimators as alternatives, and Barndorff-Nielsen--Shephard, Hayashi--Yoshida and Zhang for covariance (pp.~317--318), none read here.

\qffield{Pitfalls} The legacy model said realised variance ``can diverge as sampling becomes too fine'' and the Epps definition credited tick effects; both hold under stated noise and asynchrony, not for every series. Its remedy for covariance, ``refresh-time or previous-tick sampling'', is not in the passages read: previous-tick sampling itself produces the Epps effect in the check above. \emph{Refresh time} (read 2026-10-16). Barndorff-Nielsen, Hansen, Lunde and Shephard define it for $d$ assets with trade times $t^{(i)}_j$ and counts $N^{(i)}(t)$ as $\tau_1=\max_i t^{(i)}_1$ and $\tau_{j+1}=\max_i t^{(i)}_{N^{(i)}(\tau_j)+1}$, the first time after $\tau_j$ at which every price has been updated; they retain $p=dN/\sum_i n^{(i)}$ of the data (21/27 in their three-asset figure), the sample falls with $\log d$ when trades arrive as Poisson processes, and they treat the vector at $\tau_j$ as if fully updated although ``we only see a single new price and $d-1$ stale prices'', showing that this staleness does not affect the asymptotic distribution of their \emph{realised kernel}, which includes lagged cross-products \citep[sec.~2.1.1, Definition~1 with nn.~2--3 and Figure~1, pp.~4--5]{bnhls-2008}. Their limit theory assumes noise that is covariance stationary given the efficient price, with summable autocovariances (Assumption~1, not white noise), and refresh times $\tau_j=T(j/n)$ from a strictly positive c\`adl\`ag time change adapted to the filtration (Assumption~2, pp.~7--8). Refresh time is therefore a synchronisation step, not a covariance estimator on its own. Checked (\nolinkurl{notes/artefacts/execution-market-making/source-followup-checks.py}, section~1): two Brownian motions with correlation 0.6 observed at independent Poisson times (400 per unit time, no noise) give a mean realised covariance over 60 paths of 0.39 at refresh times, 0.60 once one lead and one lag are added, 0.59 for the Hayashi--Yoshida sum of products of overlapping returns, and 0.02 at a previous-tick grid of 1/5000. The Hayashi--Yoshida estimator is the paper's sum $U_n=\sum_{i,j}\Delta P^1(I^i)\Delta P^2(J^j)\mathbf 1_{\{I^i\cap J^j\neq\emptyset\}}$ of products of returns over overlapping observation intervals \citep[Definition~3.1, eq.~(3.1), p.~368]{hayashi-yoshida-2005}. Its consistency (Theorem~3.1, p.~369: in $L^2$ and in probability) holds for diffusions with deterministic, bounded volatilities and deterministic correlation (p.~367) under Condition~C: sampling intervals independent of the prices, and the expected longest interval going to zero (p.~368); without drift it is unbiased for every $n$ (Remark~3.3, p.~369), and realised covariance at a fixed interval is not consistent (Remark~3.4, p.~375). Microstructure noise is not in the model; the authors mention it only as a possible extension (p.~367). The simulation above meets these conditions (independent Poisson times, no noise), so its 0.59 is what the theorem predicts; the observation times of real trades depend on prices, and real prices carry noise. In Barndorff-Nielsen et al.'s data (ten NYSE stocks, TAQ, 3 January 2005 to 29 June 2007, 626 days, trades and mid-quotes), the Hayashi--Yoshida estimator was ``even more distorted'' than the 15-second realised covariance relative to open-to-close covariances, which noise-free simulations cannot show (sec.~5.3, p.~18). In pairs, refresh time ``rarely'' cost more than half the observations of the most frequently traded stock and typically 30 to 40 percent of the least active one; for all ten together more than 20 percent remained in the worst case (Table~3, pp.~17--18; page numbers are the paper's printed ones).

\qffield{See also} \hyperref[def:volatility-components]{fundamental and transitory volatility}; \hyperref[res:roll-spread-estimator]{Roll's estimator} (bounce as negative autocovariance); \hyperref[meth:garman-klass-volatility]{Garman--Klass volatility}; \hyperref[ev:high-frequency-stylised-facts]{stylised facts}.

\qffield{Further reading} \citep{bnhls-2008} for refresh time and realised kernels; \citep{hayashi-yoshida-2005} for the overlapping-interval estimator and its conditions. The legacy blocks are kept verbatim in \nolinkurl{notes/artefacts/market-simulation/}, with the simulations in \nolinkurl{order-flow-checks.py}.

\qffield{Index} realised variance; integrated variance; realised covariance; signature plot; volatility signature plot; variogram; microstructure noise; bid--ask bounce; Epps effect; asynchronous trading; previous-tick sampling

\qffield{Status} Verification: partly checked 2026-10-16 (against the text layers of local copies of \texttt{alpha-laruelle}, 2018 World Scientific printing, pp.~316--319, and \texttt{catalogue-0528}, Cambridge 2018, pp.~25--28; noise and asynchrony simulated; the limits restricted to continuous processes, the jump and co-jump limits and the noise-model assumptions added 2026-10-16, with a jump counterexample computed; source follow-up of 2026-10-16: Barndorff-Nielsen, Hansen, Lunde and Shephard, CREATES Research Paper 2008-63, secs.~2.1 and~2.3 and~5.1--5.3, PDF pp.~5--9 and 17--19, with refresh time and the Hayashi--Yoshida estimator simulated; second follow-up 2026-10-16: \texttt{hayashi-yoshida-2005}, published version, sec.~3, pp.~367--375, Condition~C and Definition~3.1 checked by page image). Stability: draft. Currency: estimators timeless; illustrations from 2012 data.
\end{qfmethod}

\begin{qfmethod}[Intraday deseasonalisation: additive and multiplicative decompositions, centred moving averages, seasonal coefficients and a factorised volume model]\label{meth:intraday-deseasonalisation}
\qffield{Statement} A series is decomposed into a trend $C_t$, a seasonal component $S_t$ and an irregular part $\varepsilon_t$, additively $X_t=C_t+S_t+\varepsilon_t$ or multiplicatively $X_t=C_tS_t+\varepsilon_t$ \citep[app.~A, sec.~A.11.1, p.~299]{alpha-laruelle}. A centred moving average of odd length $2k+1$ is $MA_t=(X_{t-k}+\dots+X_{t+k})/(2k+1)$; of even length $2k$ it weights the end points by $\tfrac12$, $MA_t=(\tfrac12X_{t-k}+X_{t-k+1}+\dots+X_{t+k-1}+\tfrac12X_{t+k})/(2k)$, and it cannot be computed for the first and last $k$ observations (pp.~299--300). For period $p$ and $t=(i-1)p+j$, with $MA$ of length $p$ estimating the trend: additively, $\hat S_j=\tfrac1n\sum_i(X_{ij}-MA_{ij})$, centred as $S_j=\hat S_j-\tfrac1p\sum_k\hat S_k$; multiplicatively, $\hat S_j=\tfrac1n\sum_iX_{ij}/MA_{ij}$, normalised as $S_j=\hat S_j/(\tfrac1p\sum_k\hat S_k)$, so that the coefficients sum to $p$; means may be replaced by medians (pp.~300--302).

\qffield{Conditions} The moving average's length must equal the period of the seasonal variation or a multiple of it, so that it carries no seasonality (p.~300). The multiplicative estimate assumes $\varepsilon_{ij}$ small relative to the trend (p.~301). Interpretation, checked (\nolinkurl{order-flow-checks.py}): a centred average uses $k$ later observations, so a shock at $t$ changes the estimates for the preceding $k$ times; seasonal coefficients estimated this way on a sample cannot be used at a decision time inside that sample without look-ahead.

\qffield{Practice} Lehalle and Laruelle's volume model for slot $n$, day $d$, month $m$ and year $y$ is $V(y,m,d,n,\omega)=(f(n)+g(d,n))F(m)G(y)v(n,\omega)$, with $f$ the mean slot profile, $g$ special-day effects such as expiries or macroeconomic news, $F$ and $G$ monthly and yearly effects and $v$ random; with $E[H]=1$ and $\sum_mF(m)=1$, $E[V_M(y,m,\cdot)]=G(y)F(m)$ and $E[V_Y(y,\cdot)]=G(y)$ (sec.~A.11.2, pp.~302--303). Their Total SA example (2 January 2007 to 1 June 2012) shows a bump at noon on Fridays from derivative expiries (pp.~303--304). Bouchaud et al.'s alternative for point processes is a \emph{business time} $\hat t(t)=\int^t\bar\phi(u)\,du$ built from the average activity profile, which flattens it \citep[chap.~9, sec.~9.3.1, p.~173]{catalogue-0528}; see \hyperref[def:market-clocks-durations]{market clocks}. Checked on a synthetic series with 13 slots and 60 days: multiplicative coefficients recovered to within 0.005.

\qffield{Pitfalls} The legacy notes wrote the even-length average as $A_t$ and asked for the window to ``match'' the period; Lehalle and Laruelle require it to equal the period or a multiple. The practitioner note ``remove predictable clock-time effects before comparing venues, estimating impact, or training intraday features'' is sound as a warning that a noon lull, a close auction, an expiry day or a news window can look like a signal, but the removal must use profiles estimated from earlier data (interpretation). \emph{Mean against path} (slice~36, from the legacy data-engineering file's \emph{Averaging Effect}). ``To show that the average profile of a quantity does not mean that this profile exists in the real world'', Lehalle and Laruelle compare Adidas's mean intraday volatility from 27 April to 27 July 2012 with its path on 25 May 2012, which ``does not have the same behavior as the mean'' \citep[app.~A, sec.~A.14.1, pp.~318--320, Figure~A.27]{alpha-laruelle}. A seasonal profile is an expectation, not a typical day: in a simulation of 500 calm days each with one short spike at a random time, the average profile peaks at 1.36 while every day peaks at 7 (\nolinkurl{slice36-checks.py}). The legacy caveat that averaging across days hides event days, closing auctions and isolated liquidity shocks follows; a simulator validated only on average profiles can miss the day-level behaviour.

\qffield{See also} \hyperref[ev:high-frequency-stylised-facts]{stylised facts} (intraday patterns); \hyperref[def:market-clocks-durations]{market clocks}; \hyperref[pit:label-weight-leakage]{time-validity of labels and fitted transforms}; \hyperref[def:execution-benchmark]{execution benchmarks} (volume curves).

\qffield{Further reading} The legacy blocks are kept verbatim in \nolinkurl{notes/artefacts/market-simulation/}. The averaging-effect block merged in slice~36 is kept in \nolinkurl{notes/artefacts/lob-deep-learning/}.

\qffield{Index} seasonality; deseasonalisation; additive model; multiplicative model; trend; centred moving average; seasonal coefficient; intraday volume profile; volume model; special days; expiry effect; business time; mean against path; averaging effect

\qffield{Status} Verification: partly checked 2026-10-16 (against the text layers of local copies of \texttt{alpha-laruelle}, 2018 World Scientific printing, pp.~298--304, and \texttt{catalogue-0528}, Cambridge 2018, p.~173; a synthetic recovery and the look-ahead computed; slice~36: pp.~318--320 for the averaging effect). Stability: draft. Currency: timeless; the illustration from 2007--2012 data.
\end{qfmethod}

\begin{qfdefinition}[Market clocks and durations: clock, event, trade, quote, volume and business time, and when each is the right index]\label{def:market-clocks-durations}
\qffield{Statement} For event times $0<\tau_1<\tau_2<\cdots$ the \emph{duration} is $\Delta\tau_i=\tau_i-\tau_{i-1}$ \citep[app.~A, sec.~A.12.1, Definition~5, p.~306]{alpha-laruelle}. An analysis can be indexed by an event counter (trades, quote revisions, or both) or by wall-clock time at fixed intervals; Hasbrouck's rule is that for a single precisely time-stamped security an event index is ``probably the first choice'', because the process is more likely to be covariance stationary in event time, while comparisons across securities may favour wall-clock time, volatility per minute being more comparable than volatility per trade \citep[sec.~9.7.3, p.~90]{s6-hasbrouck}. \emph{Business time} rescales clock time by the average activity profile $\bar\phi$, $\hat t(t)=\int^t\bar\phi(u)\,du$, so that intraday activity is flat \citep[chap.~9, sec.~9.3.1, p.~173]{catalogue-0528}. Volume and dollar clocks, one step per fixed quantity or value traded, are the sampling rules of \hyperref[meth:activity-sampled-bars]{activity-sampled bars}.

\qffield{Conditions} With a wall-clock index, several events in one interval must be aggregated or transformed (Hasbrouck, p.~90). Propagator models are almost always written in trade time and Hawkes processes in physical time; hybrids exist \citep[app.~A, sec.~A.12.3, p.~311]{alpha-laruelle}. For a homogeneous Poisson process the durations are i.i.d.\ exponential with mean $1/\lambda$ (p.~306), so high intensity means short expected durations; for clustered processes the expected duration from a given time depends on the recent history.

\qffield{Practice} Lehalle and Laruelle note that replacing physical time by event time ``would change a lot the meaning of the obtained models'' (chap.~3, p.~227). Interpretation: latency, risk limits, news and benchmark windows are defined in clock time, and an event-time or business-time model must be mapped back before it is used for them; a transformation that flattens activity also hides the periods of intense activity a risk manager cares about, so diagnostics are kept in both clocks.

\qffield{Pitfalls} The legacy definition called the gap between any events a ``trade duration''; Hasbrouck uses duration analysis for other waiting times too, such as a limit order's time to execution, whose observations are censored when the order is cancelled first \citep[sec.~15.2.3, pp.~163--164]{s6-hasbrouck}. Its ``quote time'' and ``business time'' entries are covered above; ``business time'' there meant activity-adjusted time ``after removing intraday seasonality'', which matches Bouchaud et al.'s construction only when $\bar\phi$ is estimated from data before the period it rescales.

\qffield{See also} \hyperref[def:order-event-point-processes]{point processes and intensities}; \hyperref[meth:activity-sampled-bars]{activity-sampled bars}; \hyperref[meth:intraday-deseasonalisation]{intraday deseasonalisation}; \hyperref[mod:propagator]{propagator model} (trade time).

\qffield{Further reading} Engle and Russell's autoregressive conditional duration model is named by Bacry et al.\ (\texttt{wide-bacry-hawkes}, p.~1) but not read. The legacy file is kept verbatim in \nolinkurl{notes/artefacts/market-simulation/}.

\qffield{Index} duration; inter-arrival time; event time; trade time; quote time; wall-clock time; calendar time; business time; volume time; covariance stationarity; time deformation

\qffield{Status} Verification: partly checked 2026-10-16 (against the text layers of local copies of \texttt{alpha-laruelle}, 2018 World Scientific printing, pp.~227, 306 and~311, \texttt{s6-hasbrouck}, 2007 Oxford printing, pp.~90 and 163--164, and \texttt{catalogue-0528}, Cambridge 2018, p.~173). Stability: draft. Currency: timeless.
\end{qfdefinition}

\begin{qfdefinition}[Point processes for order-book events: counting process, conditional intensity, the homogeneous Poisson baseline and the time-change theorem]\label{def:order-event-point-processes}
\qffield{Statement} A simple point process is an increasing sequence of event times $(\tau_i)_{i\ge1}$, for example trades or limit-order arrivals; its \emph{counting process} is $N_t=\sum_{i\ge1}\mathbf 1\{\tau_i\le t\}$ \citep[app.~A, sec.~A.12.1, Definitions~3--4, pp.~305--306]{alpha-laruelle}. For $N$ adapted to a filtration $\mathcal F_t$, the left-continuous \emph{intensity} is $\lambda(t\mid\mathcal F_t)=\lim_{h\downarrow0}h^{-1}E[N_{t+h}-N_t\mid\mathcal F_t]$, or equivalently $\lim_{h\downarrow0}h^{-1}P(N_{t+h}-N_t>0\mid\mathcal F_t)$; it depends on the filtration, taken to be the process's own history (Definition~6, p.~306). A homogeneous Poisson process has $P(N_{t+h}-N_t=1\mid\mathcal F_t)=\lambda h+o(h)$, an intensity that ignores the history, and i.i.d.\ exponential durations (p.~306). With the integrated intensity $\Lambda(\tau_{i-1},\tau_i)=\int_{\tau_{i-1}}^{\tau_i}\lambda(s)\,ds$, the \emph{time-change theorem} states that $\tilde N_\tau=N_{t_\tau}$, with $\int_0^{t_\tau}\lambda(s)\,ds=\tau$, is a unit-rate Poisson process (Theorem~1, pp.~307--308).

\qffield{Notation} Lehalle and Laruelle model order-book events with three intensities, $\lambda^+$ for limit-order insertions, $\lambda^c$ for cancellations and $\lambda^x$ for transactions \citep[chap.~3, p.~221]{alpha-laruelle}. Bouchaud et al.\ write $\phi(t)$ for the intensity and define the clustering ratio $r(\tau)=V[N(t+\tau)-N(t)]/E[N(t+\tau)-N(t)]$, equal to 1 for a Poisson process, above 1 for clustering and below 1 for inhibition \citep[chap.~9, secs.~9.1--9.1.2, pp.~164--165]{catalogue-0528}.

\qffield{Conditions} The intensity is conditioned on information strictly before $t$: interpretation, the intensity must be predictable (left-continuous), which is what makes it usable for prediction and simulation. Lehalle and Laruelle print the theorem's condition as $\int_0^\infty\lambda(s)\,ds=1$ (p.~308, checked by page image); the integrated intensity must instead diverge, $\int_0^\infty\lambda(s)\,ds=\infty$, for the rescaled times to cover the half-line (interpretation when written; the half-line statement, attributed to Meyer (1969) and Papangelou (1972), was then not read in a point-process text). \emph{Half-line statement} (batch of 2026-10-17). Laub, Taimre and Pollett restate Daley and Vere-Jones's Proposition~7.4.IV: for ``an unbounded, increasing sequence of time points $\{t_1,t_2,\dots\}$ in the half-line $(0,\infty)$, and a monotonic, continuous compensator $\Lambda(\cdot)$ such that $\lim_{t\to\infty}\Lambda(t)=\infty$ a.s.'', the transformed sequence $\{\Lambda(t_1),\Lambda(t_2),\dots\}$ is a realisation of a unit-rate Poisson process if and only if the original is a realisation of the point process defined by $\Lambda$ \citep[arXiv v1, July 2015, sec.~5.1, Theorem~5, p.~18]{laub-taimre-pollett-2015}; their finite-window version (Theorem~4, after Brown et al.) needs a positive intensity and $\Lambda(T)<\infty$. This confirms the divergence condition stated above and that Lehalle and Laruelle's ``$=1$'' is a misprint; Daley and Vere-Jones's own text was not read, so the attribution is through this survey. \emph{Finite-window statement} (read 2026-10-16). Brown, Barbieri, Ventura, Kass and Frank state the theorem for event times $0<u_1<\dots<u_n<T$ of a point process whose conditional intensity satisfies $0<\lambda(t\mid H_t)$ for all $t\in(0,T]$, with $\Lambda(u_k)=\int_0^{u_k}\lambda(u\mid H_u)\,du$ and $\Lambda(t)<\infty$ with probability one for all $t\in(0,T]$: then the $\Lambda(u_k)$ ``are a Poisson process with unit rate'', proved by showing that $\tau_k=\Lambda(u_k)-\Lambda(u_{k-1})$ are i.i.d.\ unit exponential, including the censored interval after $u_n$ \citep[sec.~2.2, eqs.~(2.6)--(2.11), p.~329]{brown-time-rescaling-2002}. Existence of a conditional intensity makes the compensator absolutely continuous, so the continuity condition is built in; on a finite window no divergence is needed, since only $(0,\Lambda(T)]$ is covered. Applied as a goodness-of-fit test, a correctly specified model's rescaled durations are i.i.d.\ unit exponential, and $z_k=1-e^{-\tau_k}$ are i.i.d.\ uniform; Brown et al.\ plot the ordered $z_{(k)}$ against $b_k=(k-\tfrac12)/n$ with 95 percent Kolmogorov--Smirnov bounds $b_k\pm1.36/\sqrt n$ for moderate to large $n$ (eqs.~(2.16)--(2.17), p.~331). Those bounds are for a fully specified intensity; with parameters estimated from the same data they are only approximate (interpretation, not a statement of the paper).

\qffield{Practice} Checked (\nolinkurl{order-flow-checks.py}): durations of a simulated exponential-kernel Hawkes process are far from exponential (Kolmogorov--Smirnov $p<10^{-260}$), while the same durations rescaled by the true compensator have mean 1.005 and are not rejected as unit exponential ($p=0.35$). The Poisson baseline is diagnostic: as \hyperref[mod:queue-reactive-intensities]{queue-reactive intensities} shows, context-free intensities do not reproduce queue sizes.

\qffield{Pitfalls} The legacy definitions wrote $\mathcal F_t$ as ``the information available before $t$'' without the left-continuity that makes it so, and stated the time change without its conditions. The legacy statement that a Poisson baseline ``cannot reproduce realistic queue-size distributions because event rates are strongly state-dependent'' is Lehalle and Laruelle's finding for one large-tick stock (Figure~3.11), not a general theorem.

\qffield{See also} \hyperref[mod:hawkes-order-flow]{Hawkes processes for order flow}; \hyperref[mod:queue-reactive-intensities]{queue-reactive intensities}; \hyperref[def:market-clocks-durations]{market clocks and durations}.

\qffield{Further reading} Daley and Vere-Jones, cited by Lehalle and Laruelle for point processes, is not read (its Proposition~7.4.IV is cited through Laub, Taimre and Pollett's survey); Brown et al.\ (2002) was read for the finite-window theorem and its goodness-of-fit use. The legacy file is kept verbatim in \nolinkurl{notes/artefacts/market-making-lob/}.

\qffield{Index} point process; counting process; conditional intensity; stochastic intensity; filtration; homogeneous Poisson process; exponential durations; integrated intensity; compensator; time-change theorem; time rescaling; clustering ratio; insertion intensity; cancellation intensity; transaction intensity

\qffield{Status} Verification: partly checked 2026-10-16 (against the text layer of a local copy of \texttt{alpha-laruelle}, 2018 World Scientific printing, pp.~221 and 305--308, page images of pp.~306 and~308; and \texttt{catalogue-0528}, Cambridge 2018, pp.~164--165; the time change computed; source follow-up of 2026-10-16: Brown et al., \emph{Neural Computation} 14 (2002), from the copy on a co-author's Carnegie Mellon page, pp.~325--331; batch of 2026-10-17: Laub, Taimre and Pollett, arXiv v1, p.~18, text layer). Stability: draft. Currency: timeless.
\end{qfdefinition}

\begin{qfmodel}[Hawkes processes for order flow: the linear self-exciting intensity, its conditions, branching ratio and stationarity, multivariate and marked forms, inhibition through a non-linear link, and what calibrations show]\label{mod:hawkes-order-flow}
\qffield{Statement} A linear Hawkes process has intensity
\[ \lambda(t)=\lambda_0(t)+\int_{-\infty}^t\nu(t-s)\,dN_s=\lambda_0(t)+\sum_{\tau_i<t}\nu(t-\tau_i), \]
with a deterministic baseline $\lambda_0\ge0$ and a kernel $\nu\ge0$ that ``expresses the positive influence of the past events'' \citep[app.~A, sec.~A.12.2, Definition~7, p.~308]{alpha-laruelle}. Bacry, Mastromatteo and Muzy's multivariate definition, $\lambda^i_t=\mu^i+\sum_j\int\varphi_{ij}(t-t')\,dN^j_{t'}$, requires each kernel component to be non-negative, causal ($\varphi_{ij}(t)=0$ for $t<0$) and integrable, and the process has asymptotically stationary increments if the matrix of kernel norms $\|\Phi\|=\{\int|\varphi_{ij}|\}$ has spectral radius below 1 \citep[arXiv v2, sec.~2.1, Definition~1 and Proposition~1, pp.~3--4]{wide-bacry-hawkes}. In one dimension, with norm $g=\int_0^\infty\nu(u)\,du<1$, the mean intensity is $\bar\lambda=\lambda_0/(1-g)\ge\lambda_0$; $g$ is the \emph{branching ratio}, the expected number of events directly triggered by one event, and a family of events has mean size $1/(1-g)$ \citep[chap.~9, secs.~9.2.2 and~9.2.6, pp.~167--168 and~172]{catalogue-0528}. For $\nu(u)=\alpha e^{-\beta u}\mathbf 1\{u\ge0\}$ and constant $\lambda_0$, $g=\alpha/\beta$ and $E[\lambda(t)]=\lambda_0/(1-\alpha/\beta)$ (Lehalle and Laruelle, pp.~309 and~311).

\qffield{Notation} Parameterisations differ: Lehalle and Laruelle's $\alpha e^{-\beta t}$ has norm $\alpha/\beta$; Bacry et al.'s $\alpha\beta e^{-\beta t}$ (p.~5) and Bouchaud et al.'s $g\omega e^{-\omega t}$ (eq.~(9.21), p.~170) have norms $\alpha$ and $g$. A fitted ``$\alpha$'' must be read with its kernel.

\qffield{Conditions} \emph{Excitation only.} A non-negative kernel is necessary and sufficient for the intensity to stay non-negative for every realisation, and always gives self-excitation and clustering \citep[chap.~9, sec.~9.5, p.~181]{catalogue-0528}. \emph{Inhibition} needs a non-linear form, $\lambda(t)=F\bigl(\lambda_0+\int\nu(t-u)\,dN_u\bigr)$ with $F\ge0$, such as the positive part or $\ln(1+e^x)$; if $F$'s slope never exceeds 1 the same stability condition holds (Brémaud and Massoulié, 1996, as reported, p.~181). Zhu's thesis restates their theorems, for $\lambda_t=\lambda\bigl(\int_{-\infty}^th(t-s)N(ds)\bigr)$, in two parts \citep[arXiv v3, sec.~1.5, Theorems~4--5, pp.~33--34]{zhu-nonlinear-hawkes-2013}:
\begin{itemize}
\item \emph{Existence} of a stationary process needs $\lambda(\cdot)$ non-negative, non-decreasing and left-continuous with $\lambda(z)\le C+\alpha z$ and $\alpha\int_0^\infty h<1$, without Lipschitz continuity (Theorem~4).
\item \emph{Uniqueness} of the stationary distribution with finite intensity, and stability in distribution, need $\lambda(\cdot)$ $\alpha$-Lipschitz with $\alpha\|h\|_{L^1}<1$. Stability in variation needs in addition $\int_0^\infty t\,h(t)\,dt<\infty$ (Theorem~5).
\end{itemize}
Zhu notes that the original allows $h$ to take negative values, while his own work keeps $h\ge0$. \emph{The original, read in the third batch of 2026-10-17} (an Internet Archive capture of the Project Euclid PDF). Brémaud and Massoulié write $\lambda(t)=\phi\bigl(\int_{(-\infty,t)}h(t-s)N(ds)\bigr)$ with $\phi:\mathbb R\to\mathbb R^+$ and $h:\mathbb R^+\to\mathbb R$, so a signed kernel is allowed; the linear Hawkes process is the case $h\ge0$, $\phi(x)=\nu+x$ \citep[eq.~(1), p.~1564]{bremaud-massoulie-1996}. For a signed kernel the norm is that of $|h|$: Theorem~1 requires $\phi$ $\alpha$-Lipschitz and $\alpha\int_{\mathbb R^+}|h(t)|\,dt<1$ (eq.~(3)), and then gives a unique stationary distribution with finite intensity and stability in distribution; stability in variation needs in addition $\int t|h(t)|\,dt<\infty$ (eq.~(4), p.~1567). Theorem~2 drops the bound on the $L^1$ norm when $\phi$ is also bounded: $h\in L^1$ and~(4) then suffice (p.~1569). So ``slope never exceeds 1'' is the $\alpha=1$ case and needs $\int|h|<1$, not $\int h<1$; the ``same stability condition'' as reported holds only when the inhibiting part is counted in absolute value. Such models lose most analytical results, while simulation and calibration remain possible (Bacry et al., sec.~2.2.4, p.~8). Checked (\nolinkurl{order-flow-checks.py}): with $\lambda_0=0.5$ and a kernel $-0.8e^{-2u}$, the linear intensity just after an event is $-0.3$, not an intensity; with the positive-part link the simulated rate is 0.37 and the clustering ratio 0.62, below 1 as for repelling events. \emph{Stationarity.} With $g\ge1$ the process explodes or, at $g=1$, is critical; the stationary formulas then fail (Bouchaud et al., p.~168). \emph{Multivariate and marked flow.} Market orders can trigger market orders, limit orders and cancellations, and a $K$-type process has $\bar\phi=(I-g)^{-1}\phi_0$ with $g_{ij}=\int\Phi_{ij}$ (sec.~9.5, pp.~182--183); Lehalle and Laruelle sketch classes of agents (banks, brokers, high-frequency traders) with buy and sell intensities excited across classes under $\|\varphi\|_1<1$ \citep[app.~A, sec.~A.12.3, pp.~315--316]{alpha-laruelle}.

\qffield{Practice} \emph{Calibrations.} Lehalle and Laruelle fit a one-dimensional exponential Hawkes process to trades in France Telecom on the Chi-Delta dark pool on 23 July 2012, with $\lambda_0=140$, $\alpha=250$, $\beta=300$, a branching ratio of 0.83 (Figure~A.24, p.~310). Bouchaud et al.\ report maximum-likelihood fits across many contracts with $g$ close to 1 and power-law kernels, the short time scale falling from seconds in 1998 to milliseconds in 2012, and TSLA mid-price changes whose effective feedback approaches 1 as the window grows (pp.~174--175 and~178). Bacry et al.\ report, from Bacry and Muzy's model of EuroStoxx and Euro-Bund futures over 800 trading days in 2009--2012, slightly negative self-kernels in the feedback of price moves on market-order flow: an up move raises sell intensity and lowers buy intensity, an inhibitory effect \citep[arXiv v2, sec.~5.1, p.~27]{wide-bacry-hawkes}. Checked by simulation: an exponential-kernel process with branching ratio 0.6 has mean rate 2.47 against the theoretical 2.50 and a clustering ratio of 6.4 in windows of 50 against a long-window limit of 6.25.

\qffield{Pitfalls} The legacy taxonomy said a marked Hawkes model ``lets one class of event excite or inhibit another''; in the linear model, kernels are non-negative and only excitation is possible. Calibration ``does not prove the existence of a genuine causal self-excitation mechanism'': two-point correlations do not imply causality, and near-critical fits may mean either fragile markets or a framework too restrictive (Bouchaud et al., pp.~169, n.~4, and 184). Linear Hawkes processes do not reproduce strong volatility clustering and fat tails without extensions such as price feedback (p.~185). Time-stamp errors distort maximum-likelihood fits of power-law kernels (p.~178).

\qffield{See also} \hyperref[def:order-event-point-processes]{point processes and intensities}; \hyperref[mod:queue-reactive-intensities]{queue-reactive intensities} (Hawkes memory as a remedy); \hyperref[ev:high-frequency-stylised-facts]{stylised facts} (activity clustering); \hyperref[mod:propagator]{propagator model} (event-time alternative).

\qffield{Further reading} \texttt{wide-bacry-hawkes} read in its arXiv version 2 (17 May 2015), secs.~1--2.2.4 and the passages of secs.~5.1 and~6.1 cited; the journal version was not checked. Hawkes (1971), Brémaud and Massoulié (2001) and Bacry and Muzy are not read; Brémaud and Massoulié (1996) was read for eq.~(1) and Theorems~1--2 (third batch of 2026-10-17). The legacy file is kept verbatim in \nolinkurl{notes/artefacts/market-making-lob/}.

\qffield{Index} Hawkes process; self-exciting process; mutually exciting process; kernel; exponential kernel; power-law kernel; branching ratio; kernel norm; stationarity; spectral radius; criticality; endogeneity; multivariate Hawkes; marked point process; inhibition; non-linear Hawkes; clustering ratio; maximum likelihood

\qffield{Status} Verification: partly checked 2026-10-16 (against the text layers of local copies of \texttt{alpha-laruelle}, 2018 World Scientific printing, pp.~304--311 and 315--316, page image of p.~308; \texttt{catalogue-0528}, Cambridge 2018, pp.~163--185; and \texttt{wide-bacry-hawkes}, arXiv v2, pp.~1--8 and 27--30, downloaded 2026-10-16; the excitation, inhibition and stationarity cases simulated; batch of 2026-10-17: Brémaud and Massoulié's Theorems as restated by Zhu, arXiv:1304.7531v3, pp.~33--34, text layer; third batch of 2026-10-17: \texttt{bremaud-massoulie-1996}, pp.~1563--1564, 1567 and~1569, page images). Stability: draft. Currency: calibrations from 1998--2015 data.
\end{qfmodel}
